\documentclass[11pt,a4paper]{article}
\usepackage[utf8]{inputenc}
\usepackage[T1]{fontenc}
\usepackage[polish]{babel}
\usepackage{lmodern,amsmath,amssymb,amsthm,mathtools,booktabs,longtable,array,microtype}
\usepackage[margin=24mm]{geometry}
\usepackage{xcolor,hyperref}
\hypersetup{colorlinks=true,linkcolor=blue!45!black,urlcolor=blue!50!black,citecolor=blue!45!black}
\newtheorem{tw}{Twierdzenie}[section]
\newtheorem{lem}[tw]{Lemat}
\newtheorem{wn}[tw]{Wniosek}
\theoremstyle{definition}\newtheorem{defn}[tw]{Definicja}
\theoremstyle{remark}\newtheorem*{uw}{Uwaga}
\newcommand{\Pp}{\mathbb P}
\newcommand{\R}{\mathcal R}
\newcommand{\B}{\mathcal B}
\newcommand{\T}{\mathcal T}
\newcommand{\one}{\mathbf1}
\newcommand{\D}{\mathcal D}
\newcommand{\HH}{\mathcal H}
\newcommand{\PP}{\mathrm{PP}}
\newcommand{\Ptwo}{\mathrm{P}_2}
\newcommand{\eps}{\varepsilon}
\DeclareMathOperator{\ord}{ord}
\title{\textbf{Reguła Krzyża, rytm 4--6--0 i bariera parzystości}\\[4pt]\large Wspólny bilans CRT--Möbiusa, ruchomy Czybyszew\\i dodatnie jądra Newtona\\\normalsize Audyt etapów 191--209 oraz synteza wcześniejszego programu v10}
\author{Adam Berger --- projekt badawczy\\\small Opracowanie i audyt z pomocą ChatGPT (GPT-6.2)}
\date{26 listopada 2026 r.\quad Wersja audytowa 1.0}
\begin{document}
\maketitle
\begin{abstract}
Przedstawiamy jeden spójny zapis konstrukcji dwóch ramion, rytmu 4--6--0,
proporcji CRT $25/17$, momentu Liouville'a, naturalnych okien Czybyszewa,
oraz podpisanej reszty Möbiusa B1. Dodatnie wielomiany Newtona pozwalają
wygasić wkład dowolnej skończonej liczby bliźniaków. Dla $a=3730$, $r=5120$,
$u=1/20$ odtworzono wszystkie $10^8$ współczynników i uzyskano skierowaną
arytmetyką przedziałową certyfikat $10^{2619}\B_{a,r}(u)<-20{,}49$, obejmujący
również nieskończony ogon. Wyprowadzamy także wspólne kryterium B1 z dodatnim
jądrem na blokach dyadycznych. Żaden wynik skończony nie rozstrzyga jednak
znaku w granicy $u\downarrow0$. Dokument nie przedstawia dowodu nieskończoności
bliźniaków: precyzuje brakujący lemat analityczny. Plik Lean formalizuje
rdzeń algebraiczny i logiczny redukcji, nie zastępując otwartych oszacowań.
\end{abstract}
\clearpage
\tableofcontents
\clearpage
\section{Status i zakres audytu}
W całej pracy rozróżniamy: (i) tożsamość ścisłą, (ii) wniosek z klasycznego
twierdzenia analitycznego, (iii) wynik obliczeniowy, (iv) model heurystyczny,
(v) otwarty lemat. W szczególności nie utożsamiamy rozkładu kandydatów CRT z
rozkładem pierwszych; nie utożsamiamy lokalnej wariancji z globalnym błędem;
nie przyjmujemy RH, GRH, EH ani Hardy'ego--Littlewooda w redukcji B1.

Zakres zasadniczy: etapy 191--209. Włączono fundamenty v10, rytm 48 i użyteczne
wnioski etapów 178--190. Nie reprodukujemy wszystkich wcześniejszych prób
biliniowych jako dowodów: ich pełne oszacowania nie zostały niezależnie
zweryfikowane i nie są przesłankami twierdzeń poniżej. W szczególności nie
zakładamy przełamania parzystości w nowych preprintach.

\begin{center}
\begin{tabular}{p{0.56\linewidth}p{0.33\linewidth}}
\toprule
Element & Status\\\midrule
Reguła Krzyża, rytm, teleskop, momenty, jądra & tożsamości\\
Proporcja $25/17$ w pełnych okresach & ścisła\\
Jedna pierwsza i płytkie sito & BV + lemat fundamentalny\\
Dodatnie jądro i zanik skończonego wkładu & ścisłe\\
Odtworzenie danych $10^8$, znak przy $u=1/20$ & certyfikat obliczeniowy\\
Wariancja, oś $1/4$, asymptotyka $L$ & modele / pomiary\\
B1, B1 wygładzone, znak gdy $u\to0$ & otwarte\\\bottomrule
\end{tabular}
\end{center}

\section{Fundament: dwa ramiona i Reguła Krzyża}
Niech $m\ge1$ i
\[
 L_m=6m-1,\qquad C_m=12m^2-1,\qquad R_m=6m+1.
\]
Mamy dokładnie
\begin{equation}\label{eq:cross}
 L_m+R_m=12m,\qquad R_m-L_m=2,\qquad L_mR_m=3C_m+2.
\end{equation}
Ponieważ $\gcd(L_m,R_m)=1$ i oba są większe od 1, specjalny prawy stan
$\Omega(L_mR_m)=2$ jest równoważny pierwszości obu ramion. Z kolei
$\Omega(3C_m)=2$ jest równoważne pierwszości $C_m$; jednoczesny stan
specjalny $\mathrm{P}_2\mathrm{P}_2$ pary $(3C_m,L_mR_m)$ jest równoważny
$\mathrm{PPP}$. Nie mylić tego ze zwykłym stanem dwóch semipierwszych ramion.
Środkowy $C_m$ nie jest potrzebny do samego twierdzenia o bliźniakach.

Dla $p>3$ klasy ramion to $m\equiv\pm6^{-1}\pmod p$. Dla środka
$(6m)^2\equiv3\pmod p$, podczas gdy dla ramion $(6m)^2\equiv1\pmod p$.
Te zbiory są rozłączne. Dla $p>3$ i $p\nmid k$ wspólna podzielność
$C_m,C_{m+k}$ wymusza $p\mid3k^2-1$; dla $k=8$ (rytm 48) jest to $p=191$.
Wynika to z $C_{m+k}-C_m=12k(2m+k)$.

\subsection{Teleskop sita armowego v10}
Niech $I$ będzie skończonym oknem indeksów, $H=|I|$ i $n_t$ liczbą
pozycji, które przetrwały eliminatory do $t$. Przy porządku rosnącym
(właśnie dla tej wersji teleskopu), $\ell_p=n_{p^-}-n_p$,
\[
 E_p=\ell_p-\frac2p n_{p^-},\qquad M_t=H\prod_{5\le p\le t}\left(1-\frac2p\right).
\]
Stąd $n_p=(1-2/p)n_{p^-}-E_p$, $M_p=(1-2/p)M_{p^-}$ i
\begin{equation}\label{eq:tel}
 \frac{E_p}{M_p}=\frac{n_{p^-}}{M_{p^-}}-\frac{n_p}{M_p},\qquad
 \sum_{47<p\le z}\frac{E_p}{M_p}=\frac{n_{47}}{M_{47}}-\frac{n_z}{M_z}.
\end{equation}
Jeżeli $z\ge\sqrt{\max R_m}$, $n_z$ jest liczbą bliźniaków pod warunkiem,
że nie skreślamy pierwszej równej eliminatorowi (np. $\min L_m>z$).
Reguła pierwszeństwa większego eliminatora zachowuje końcowy zbiór, ale
zmienia poszczególne $E_p$; nie wolno wstawiać skreśleń liczonych malejąco do
wzoru \eqref{eq:tel} z rosnącym $p^-$. Teleskop sam nie daje dolnego oszacowania.

\section{Rytm 4--6--0 i wspólne momenty}
Eliminator 5 pozostawia $m\equiv0,2,3\pmod5$. Zatem dla par większych od
$(5,7)$ sumy ramion są $0,24,36\pmod{60}$, a końcówki mają cykl $4,6,0$.
Wprowadzamy jeden cykl
\begin{equation}\label{eq:cycle}
 T_4=(30k+11,30k+13),\quad T_6=(30k+17,30k+19),\quad
 T_0=(30k+29,30k+31).
\end{equation}
Parametry $m$ to $5k+2,5k+3,5k+5$. Kolejne przyrosty sum to $12,24,24$;
dwa przyrosty 24 dają 48 w sumie, czyli 24 w ramieniu (nie 48).
Symetria $m\mapsto-m\pmod5$ sprzęga tory 4 i 6. W szczególności
$C_4(k)=300k^2+240k+47$, $C_6(k)=300k^2+360k+107$,
$C_0(k)=300k^2+600k+299$ i $C_4(-k-1)=C_6(k)$.

\subsection{Eliminatory i dokładny okres (etap 192)}
\begin{center}
\begin{tabular}{rrrr}
\toprule $p$&klasy ramion $m$&klasy środka $m$&liczba\\\midrule
5&1,4&--&2\\7&1,6&--&2\\11&2,9&1,10&4\\13&2,11&5,8&4\\17&3,14&--&2\\19&3,16&--&2\\\bottomrule
\end{tabular}
\end{center}
Pełny okres $1616615$; w każdym torze $323323$ pozycji. Eliminacja malejąco,
z jednoznacznym przypisaniem największemu eliminatorowi z użytego zbioru:
\begin{center}
\begin{tabular}{rrr}
\toprule $p$&tylko ramiona&całe LCR\\\midrule
19&34034&34034\\17&34034&34034\\13&39270&78540\\11&39270&64260\\7&50490&32130\\\midrule
Pozostaje&126225&80325\\\bottomrule
\end{tabular}
\end{center}
Dokładna proporcja ocalałych LCR do LR wynosi $80325/126225=7/11$;
$4/11$ to wykryte złożone środki wśród ocalałych ramion. Nadwyżka $3/11$
dotyczy wyłącznie małego sita, nie stanu rzeczywistego PPP. Audyt
odtworzył także pełny okres niezależnym sitem i testem SymPy środka:
\begin{center}
\begin{tabular}{rrrr}
\toprule Tor&PP&PPP&PCP\\\midrule
4&19255&2589&16666\\6&19179&2489&16690\\0&19017&2457&16560\\
Razem&57451&7535&49916\\\bottomrule
\end{tabular}
\end{center}

\subsection{Trzeci moment i proporcja $25/17$}
Dla wskaźników dowolnego rodzaju (przesiane albo pierwsze) niech
$x_4,x_6,x_0\in\{0,1,2\}$ i $K=x_4+x_6+x_0$. W jednym cyklu
\begin{equation}\label{eq:D3}
 S_3=\binom K3,\quad H_3=x_4x_6x_0,\quad
 D_3=S_3-H_3=\sum_{i\ne j}\binom{x_i}2x_j.
\end{equation}
Dowód: trójka wybiera albo po jednym końcu z każdego toru, albo oba końce
jednego toru i jeden z pozostałych. Gdy $D_3>0$ dla wskaźników pierwszości,
występują bliźniaki. $H_3>0$ nie jest konieczne do tego wniosku.
Niech $\mathcal H=\{11,13,17,19,29,31\}$. Rodzina $\D$ ma 12 trójek
z pełnym torem, rodzina porównawcza $\mathcal J$ ma 8 trójek.
Dla $p\ge11$ trzy reszty każdej trójki są różne. Przy $p=7$ jedyna
kolizja to $17\equiv31\pmod7$. W każdej rodzinie są po dwie trójki
z tą kolizją. Współczynnik względny to $(7-2)/(7-3)=5/4$.
\begin{tw}[Dokładna proporcja CRT]\label{th:ratio}
Po dowolnym skończonym sicie na liczbach $p>5$ zawierającym 7, w pełnych
okresach CRT zachodzi
\[
 \frac{D_3^{\mathrm{sito}}}{H_3^{\mathrm{sito}}}
 =\frac{10+2\cdot5/4}{6+2\cdot5/4}=\frac{25}{17},\qquad
 \frac{S_3^{\mathrm{sito}}}{H_3^{\mathrm{sito}}}=\frac{42}{17}.
\]
\end{tw}
\begin{proof} Niezależność lokalnych warunków na pełnym okresie CRT pozwala
pomnożyć ich gęstości. Wszystkie czynniki poza 7 skracają się w ilorazie.
\end{proof}
Dla rzeczywistych trójek warunek $D_3>H_3$ odpowiada względnej skuteczności
$\rho_D/\rho_H>17/25=0{,}68$ na pełnym okresie, o ile $\rho_H>0$.
Nie mamy dolnego ograniczenia tej względnej skuteczności.

\subsection{Jedna pierwsza i płytkie sito; rytm 48}
W każdej trójce wyróżniamy jeden koniec $a$: dla $\D$ lewy koniec pełnej
pary, dla $\mathcal J$ najmniejszy koniec. Warunkowa gęstość eliminacji
$p>5$ wynosi $(\nu_p(J)-1)/(p-1)$. Przy 7 stosunek wynosi ponownie $5/4$.
Klasyczny lemat fundamentalny sita wraz z Bombierim--Winogradowem daje
asymptotyczny stosunek $25/17$ w blokach $[A,2A)$ przy
$z\to\infty$, $\log z=o(\log A)$: stosujemy sito do pierwszych w stałej
progresji modulo 30 z poziomem $D=A^{\theta}$, $0<\theta<1/2$;
$\log D/\log z\to\infty$. Błędy CRT sumuje się z wagami sita, a
pozostałe dwa końce mają wymiar 2. Jest to zastosowanie standardowych
narzędzi, nie dowód na głębokości $A^{1/3}$ (por. \cite{taosieve}).

Wyniki Tao--Teräväinena o korelacjach rzędu 1,2,3 \cite{tt} dostarczają narzędzia do badania analogicznych proporcji ze znakami Liouville'a. Przy ustalonym $z$ potrzebna jest wersja dla skończonego zbioru ustalonych progresji (jej pełną adaptację pozostawiamy poza tym audytem). Dla ustalonego skończonego zbioru klas można wybrać wspólne dobre
skale; diagonalizacja pozwala zwiększać $z$ dowolnie wolno. Nie daje jednak
jednostajności potęgowej. To samo ograniczenie występuje w rytmie 48,
przy $E=\{0,2,48,50\}$: dla $x,y\in\{0,1,2\}$
$d=xy-\one_{x>0}\one_{y>0}$ spełnia $3d\ge(x+y)(x+y-2)$.
Znaki po płytkim sicie nie certyfikują pierwszości.

W progresji $k=77t+3$ tor 0 jest zawsze złożony (dla $t\ge0$):
\[
 30k+29=7(330t+17),\qquad30k+31=11(210t+11).
\]
Trzy pierwsze wśród czterech pozostałych ramion wystarczą. Kontrprzykład
przy $k=19253$: $577601$ jest pierwsza, lecz
$577603=13\cdot157\cdot283$, $577607=37\cdot67\cdot233$,
$577609=17\cdot61\cdot557$ mają też $\lambda=-1$.

\section{Podpisany moment: precyzyjny detektor parzystości}
Odtąd $k\ge1$, $n_h(k)=30k+h$, $z(k)=\lfloor(30k+31)^{1/3}\rfloor$,
\[
 R_h(k)=\one_{P^-(n_h(k))>z(k)},\qquad
 \lambda_h(k)=(-1)^{\Omega(n_h(k))}.
\]
Wszystkie ocalałe liczby mają co najwyżej dwa czynniki pierwsze liczone
z krotnością, gdyż każdy czynnik większy od $z(k)$ jest większy od
rzeczywistego pierwiastka sześciennego. Wszystkie pierwsze ramiona są
większe od $z(k)$. Zatem dla $k\ge1$ dokładnie
\[
 R_h(k)\frac{1-\lambda_h(k)}2=\one_{n_h(k)\in\Pp}.
\]
Wariant blokowy dla $A\le k<2A$ ma wspólny próg
$z=\lfloor(60A+31)^{1/3}\rfloor$. Nie jest to ta sama funkcja co wariant
punktowy; dane etapów 197--203 są blokowe, etap 209 jest punktowy.

Dla trójki $(a,b,c)\in\D$, gdzie $(a,b)$ to pełny tor, definiujemy
\begin{align*}
 M_I&=\sum_{k\in I}\sum_{(a,b,c)\in\D}R_aR_bR_c,\\
 L_I&=\sum_{k\in I}\ell(k),\quad
 \ell(k)=\sum_{(a,b,c)\in\D}R_aR_bR_c(\lambda_a+\lambda_b),\\
 C_I&=\sum_{k\in I}\sum_{(a,b,c)\in\D}R_aR_bR_c\lambda_a\lambda_b,\\
 J_I&=\sum_{k\in I}\sum_{(a,b,c)\in\D}\one_{n_a,n_b\in\Pp}R_c.
\end{align*}
\begin{lem}[Tożsamość momentu i detekcja]\label{lem:sign}
$4J_I=M_I-L_I+C_I$, $C_I\ge-M_I$, więc $J_I\ge-L_I/4$.
Jeżeli $T(k)$ liczy pełne pierwsze tory w cyklu, to
\[
 |\ell(k)|\le24,\qquad \ell(k)\ge-8T(k).
\]
W szczególności $\ell(k)<0$ wymusza bliźniaki.
\end{lem}
\begin{proof}
Rozwinięcie $(1-\lambda_a)(1-\lambda_b)$ daje tożsamość. Na przesianym
pełnym torze suma znaków wynosi $-2,0,+2$ odpowiednio dla PP, stanów mieszanych
i $\Ptwo\Ptwo$. Każdy tor ma najwyżej cztery dopuszczalne kotwice, zatem
jeden pełny pierwszy tor wnosi co najmniej $-8$. Jest 12 trójek, każda o
wartości bezwzględnej nie większej od 2.
\end{proof}
W szczególności $T(I)\ge-L_I/8$; $J_I$ i $D_3$ nie są tym samym:
$J_I$ dopuszcza kotwicę semipierwszą, $D_3$ wymaga jej pierwszości.
Skończenie wiele bliźniaków implikuje ostateczną nieujemność $\ell$.

\section{Wyniki obliczeniowe i audyt danych}
\subsection{Ślepy test na skali $10^{33}$ (193--195)}
Początek ramion $1671875017670727170815930822929071$, milion cykli.
Zapisano prognozę przed testem pierwszości. Plik \texttt{twins.csv} zawiera
6877 par i 150 środków oznaczonych jako prawdopodobnie pierwsze.
\begin{center}
\begin{tabular}{rrrrr}
\toprule Tor&Kandydaci $z=99991$&Prognoza PP&PP&PPP\\\midrule
4&31322&2251,82&2266&49\\6&31403&2257,64&2354&41\\0&31420&2258,87&2257&60\\\midrule
Razem&94145&6768,33&6877&150\\\bottomrule
\end{tabular}
\end{center}
Audyt ponownie sprawdził wszystkie zapisane pary i wszystkie wartości $C$ przez SymPy; wyniki są zgodne. Testy GMP/SymPy były probabilistyczne na tej skali, nie formalne certyfikaty pierwszości. Na podstawie zapisanych obsad: $S_1=293744$, $S_2=34379$,
$S_3=2064$, $H_3=847$, $D_3=1217$, $D_3-H_3=370$.
Kontrmodel z zachowaniem masek i nierozdzielonych momentów, ale bez pełnego
toru, pokazuje niewystarczalność tych momentów. Nie jest modelem rzeczywistej
pierwszości. Certyfikat kontrmodelu jest w archiwum danych; audyt odtworzył wszystkie 64 częstości masek sita i ponownie sprawdził podzbiór każdej maski, brak par i trzy momenty.

\subsection{Dokładne bloki $10^3$--$10^8$ (197--199)}
\begin{center}
\begin{tabular}{rrrrr}
\toprule $A$&$M$&$L$&$D_3$&$H_3$\\\midrule
$10^3$&1812&$-874$&421&302\\
$10^4$&10197&$-5046$&2441&1627\\
$10^5$&65453&$-29664$&15031&10258\\
$10^6$&444020&$-186108$&97922&66550\\
$10^7$&3118626&$-1270646$&680450&462365\\
$10^8$&22661634&$-9074712$&4895178&3329391\\\bottomrule
\end{tabular}
\end{center}
Dla $A=10^7$ porównanie ślepego sita i pełnego sita pierwszości:
\begin{center}
\begin{tabular}{rrrr}
\toprule Próg&$D_3$&$H_3$&$D_3-H_3$\\\midrule
Bez sita&120000000&80000000&40000000\\
19&27712204&18844286&8867918\\
97&9388096&6383623&3004473\\
839&3118626&2120562&998064\\
Pierwsze&680450&462365&218085\\\bottomrule
\end{tabular}
\end{center}
Rzeczywiste pary: $332897,333289,333635$ (razem $999821$).
Tożsamość $4\cdot1129213=3118626+1270646+127580$ sprawdzona.
W skali $10^8$ proporcja rzeczywista $4895178/3329391\simeq1{,}470292$.
Dane z $10^7$ i $10^8$ nie są jednak tym samym zakresem co zapis punktowy
etapu 209.

\subsection{Rozrzedzenie i fluktuacje (200--203)}
W czterech przedziałach po $2\cdot10^8$ cykli:
\begin{center}
\begin{tabular}{rrrrr}
\toprule $A$&$T$&$M$&$L$&cykli/PP\\\midrule
$10^9$&13500520&36140835&$-13467164$&14,814\\
$10^{10}$&11333570&28011021&$-10259564$&17,647\\
$10^{11}$&9594707&21901297&$-7968534$&20,845\\
$10^{12}$&8226756&17391043&$-6319490$&24,311\\\bottomrule
\end{tabular}
\end{center}
Dla długości $w$ wprowadzono $N_{\mathrm{eff}}=w\rho_2(A)$, gdzie
$\rho_2(A)=T/B$ jest gęstością empiryczną; model bez znajomości pierwszych
to $\rho_0(A)=60C_2/\log^2(30A)$ z $C_2=\prod_{p>2}(1-(p-1)^{-2})$.
Wykładnik dopasowania $\sigma\sim N_{\rm eff}^{2\alpha}$ dla ruchomych
okien 1--10 mln w etapie 201 wyniósł $0{,}250129$, dla 1--20 mln
$0{,}244987$. Dla równej oczekiwanej masy $0{,}244293$.
Dopasowania w etapie 202 dały też przekroczenia $1/4$, wspólnie $0{,}263520$.
Nie wolno selekcjonować jedynie pomiarów bliskich $1/4$.

\section{Ruchomy Czybyszew i trzy różne osie}
Aby zapis był samodzielny, definiujemy sześciotorowy geometryczny iloczyn
(wersję znormalizowaną ruchomego dwumianu)
\[
 B_6(A,h)=\left[\prod_{c\in\HH}\prod_{j=0}^{h-1}
 \frac{30(A+h+j)+c}{30(A+j)+c}\right]^{1/6}.
\]
Całkując $1/t$ w każdym czynniku, otrzymujemy (również w nieco luźniejszej
postaci używanej wcześniej)
\begin{equation}\label{eq:cheb}
 \frac{h^2}{A+2h+109/30}\le\log B_6(A,h)\le\frac{h^2}{A+11/30}.
\end{equation}
Istotnie każdy logarytm ma postać $\int_0^h (A+j+c/30+t)^{-1}\,dt$;
wszystkie mianowniki mieszczą się między krańcami w \eqref{eq:cheb}.
Dla $h=\lfloor\sqrt A\rfloor$ mamy $B_6\to e<4$. Dla $h\le\sqrt{A\log4}$
również $B_6<4$. To czysta geometria, nie twierdzenie o pierwszych.
Przewidywana pojemność
$N_0(A,h)=60C_2 h/\log^2(30A)$ rośnie do nieskończoności dla
$h\sim\sqrt A$; rzeczywista pojemność nie jest tym samym.
\begin{center}
\begin{tabular}{rrrrr}
\toprule $A$&$h$&Średnia PP&$L/M$&$\sigma/\sqrt{N_0}$\\\midrule
$10^9$&31622&2149,19&$-0{,}365057$&2,8783\\
$10^{10}$&100000&5666,68&$-0{,}364788$&2,9574\\
$10^{11}$&316227&15173,39&$-0{,}364083$&3,1249\\
$10^{12}$&1000000&41134,13&$-0{,}363256$&2,2745\\
$10^{13}$&3162277&112742,80&$-0{,}364070$&1,9580\\\bottomrule
\end{tabular}
\end{center}
Etap 203: 780/780 ujemnych okien, 450349440 cykli. Wspólne dopasowanie
$\alpha=0{,}197577$ na $Q_{\rm eff}=N_0^2$. Dla $h\sim\sqrt A$ hipoteza
$\sigma\asymp\sqrt{N_0}$ oznacza $A^{1/4}/\log A$, czyli $Q_{\rm eff}^{1/4}$
oraz wykładnik $1/8$ względem globalnego $Q\asymp A^2$. Gdy badamy cały blok
$h\asymp A$, model reszty pierwiastkowej daje względem $Q$ wykładnik $1/4$.
To trzy różne zmienne odniesienia. Dopasowanie z danych lokalnych nie dowodzi
przedłużenia analitycznego ani położenia osobliwości. Model stałej 950 z
etapu 198 ma błąd systematyczny (w $10^8$ reszta około $67415$), więc nie
jest ustalonym członem głównym z błędem pierwiastkowym.

\section{Wersja B1: dokładny podpisany ogon Möbiusa}\label{sec:B1}
Tu pracujemy na zmiennej $n$ (lewe ramię), nie na indeksie $k$.
Niech $X<n\le2X$, $a(n)=\log n$ dla $n\in\Pp$, $n\equiv5\pmod6$,
a poza tym $a(n)=0$. Ustalmy $y=X^{0{,}49}$. Zdefiniujmy
\[
 \T(X)=\sum_{X<n\le2X}a(n)\mu^2(n+2)\Lambda(n+2),\qquad
 J_\mu(X)=\sum_{X<n\le2X}a(n)\mu(n+2).
\]
$\T$ jest ważoną liczbą bliźniaków. Nie mylić $J_\mu$ z $J_I$ z
lematu \ref{lem:sign}. Korzystając z tożsamości
\[\Lambda(N)=-\sum_{d\mid N}\mu(d)\log d,\]
otrzymujemy $\T=S_{\le y}+S_{>y}$. Po zamianie $e=(n+2)/d$ w dużym ogonie:
\begin{align}
 S_{\le y}&=-\sum_n a(n)\mu^2(n+2)\sum_{d\mid n+2,\ d\le y}\mu(d)\log d,\nonumber\\
 S_{>y}&=\sum_{\substack{e<(2X+2)/y\\(e,6)=1}}\mu(e)
 \sum_{\substack{X<n\le2X\\n+2>ey\\e\mid n+2}}
 a(n)\mu(n+2)\log\frac e{n+2}.\label{eq:large}
\end{align}
Użyto $\mu^2(N)\mu(d)=\mu(N)\mu(N/d)$ dla $d\mid N$; gdy $N$ nie jest
bezkwadratowa, obie strony są zerowe. Nie usuwamy ruchomej granicy $ey$.

\subsection{Człon mały: adaptacja lematu Murty'ego--Vatwani}
\begin{lem}\label{lem:small}
Bezwarunkowo $S_{\le y}=2C_2X+o(X)$.
\end{lem}
\begin{proof}[Uzasadnienie analityczne]
Stosujemy argument lematu 3.4 Murty'ego--Vatwani \cite{mv} z poziomem
$X^{0{,}49}$, zastępując ich warunek dystrybucyjny przez klasyczne BV.
Dla przejrzystości podajemy dostosowanie stałych. Rozwinąć
$\mu^2(n+2)=\sum_{b^2\mid n+2}\mu(b)$. Dla $b\le(\log X)^B$ moduł
$6[d,b^2]\ll X^{0{,}49}(\log X)^{2B}$ pozostaje w zakresie BV.
Waga $\log d$ i krotność po krótkim $b$ kosztują tylko potęgi logarytmu:
przy ustalonym $b$ i module możliwe $d$ dzielące go różnią się wyłącznie
dzielnikami $b$, a ich liczba jest $\le b$. Przy $b>Z=(\log X)^B$,
$d,b$ bezkwadratowe, oznaczając $g=(d,b)$, $d=gd'$, $b=gb'$,
$[d,b^2]=d'b'^2g^2$, błąd bezwzględny wynosi
\[
 \ll X\log^2X\sum_{gb'>Z}\frac{1}{g^2b'^2}
       \sum_{d'\le 3X}\frac1{d'}\ll X\frac{\log^4X}{Z}.
\]
W nierówności użyto $X/q+1\ll X/q$ dla $q\le3X$; większy moduł nie daje
żadnego dzielnika $n+2$. Analogiczny ogon członu głównego z $1/\varphi(q)$
jest zaniedbywalny po powiększeniu $B$, stosując $q/\varphi(q)\ll\log\log(q+3)$
i zbieżne sumy dla $g,b'$. Potęgi pierwszych można usunąć z BV ze stratą
$o(X)$ (liczba potęg $O(\sqrt X\log X)$, wagi dzielnikowe $X^{o(1)}$).

Dla $p\ge5$ lokalna bezkwadratowość na pierwszych w klasie 5 modulo 6 ma
czynnik $A_p=1-1/[p(p-1)]$. Jeśli $p\mid d$, warunek $p\mid n+2$ oraz
$p^2\nmid n+2$ ma gęstość $1/(p-1)-1/[p(p-1)]=1/p$.
Niech $A_* =\prod_{p\ge5} A_p$, $g_p=1/(pA_p)$. Odpowiedni szereg
$\prod_{p\ge5}(1-g_p p^{-s})=H(s)/\zeta(1+s)$ ma w $s=0$ pochodną
$H(0)$. Zbieżność ważonych sum Möbiusa przy $s=0$ wynika z PNT
(lub argumentu tauberowskiego w \cite{mv}). Stała członu małego wynosi
\[
 \frac12 A_* H(0)=\frac12\cdot3\prod_{p\ge5}
 \frac{A_p-1/p}{1-1/p}=
 \frac32\prod_{p\ge5}\left(1-\frac1{(p-1)^2}\right)=2C_2.
\]
Współczynnik $1/2$ pochodzi z progresji modulo 6, a czynnik 3 z usunięcia
pierwszych 2 i 3 w $\zeta(1+s)^{-1}$.
\end{proof}

\subsection{Człon duży i reszta podpisana}
Wprowadźmy $\Delta_e(n)=\one_{e\mid n+2}-1/\varphi(e)$ i
\begin{equation}\label{eq:R}
 \R(X)=\sum_{\substack{e<(2X+2)/y\\(e,6)=1}}\mu(e)
 \sum_{\substack{X<n\le2X\\n+2>ey}}a(n)\mu(n+2)\Delta_e(n)
       \log\frac e{n+2}.
\end{equation}
Na $n=6m-1$ jest to dokładnie $\one_{m\equiv-6^{-1}\ (e)}-1/\varphi(e)$,
nie niezważony teleskop \eqref{eq:tel}. Funkcja główna szeregu
$\sum_{(e,6)=1}\mu(e)/(\varphi(e)e^s)=H_0(s)/\zeta(1+s)$ ma
$H_0(0)=4C_2$. Standardowe oszacowania PNT dla sum Möbiusa (splot z
absolutnie zbieżnym szeregiem $H_0$) dają dla każdego ustalonego $B>0$
\[
 \sum_{e\le E,(e,6)=1}\frac{\mu(e)}{\varphi(e)}=O_B((\log E)^{-B}),\qquad
 \sum_{e\le E,(e,6)=1}\frac{\mu(e)\log e}{\varphi(e)}=-4C_2+o(1).
\]
Pierwsze oszacowanie kompensuje mnożnik $\log(n+2)$. Z $E=(n+2)/y\asymp
X^{0{,}51}$, jednostajnie po $n\in(X,2X]$, dostajemy
$S_{>y}=-4C_2J_\mu(X)+\R(X)+o(X)$. Zatem
\begin{equation}\label{eq:B1exact}
 \boxed{\T(X)=4C_2\bigl(X/2-J_\mu(X)\bigr)+\R(X)+o(X).}
\end{equation}
Dla $e<X$ oraz $n>X$ pierwszej, $(n,e)=1$; dla $(e,6)=1$ również $(-2,e)=1$.
Ortogonalność charakterów wyraża $\Delta_e$ jako
$\varphi(e)^{-1}\sum_{\chi\ne\chi_0}\overline{\chi(-2)}\chi(n)$.
Nie daje to oszacowania sumy przesuniętego Möbiusa.

\subsection{Rezerwa dodatnia i B1}
Ponieważ $\mu(n+2)\le1$, a dla $25\mid n+2$ jest zerem, PNT w jednej
zredukowanej klasie modulo 150 daje
$J_\mu(X)\le X/2-X/40+o(X)$. Stąd
\begin{equation}\label{eq:reserve}
 \T(X)\ge\delta_0 X+\R(X)+o(X),\qquad\delta_0=C_2/10>0.
\end{equation}
\begin{tw}[Redukcja B1]\label{th:B1}
Jeżeli dla nieograniczonego ciągu $X$ zachodzi
\[
 \boxed{\R(X)\ge-(C_2/20)X,}
\]
to $\T(X)>0$ dla dostatecznie dużych wyrazów tego ciągu i istnieje
nieskończenie wiele bliźniaków. Co więcej, na tych skalach ich liczba jest
$\gg X/\log^2 X$.
\end{tw}
\begin{proof} Z \eqref{eq:reserve} pozostaje $C_2X/20+o(X)$. Każda para
wnosi co najwyżej $\log^2(2X+2)$. Nieograniczone skale wymuszają nieskończoność.
\end{proof}
Dołączenie kwadratów $7^2,11^2$ zwiększa rezerwę do
$\delta_*=2C_2(1-(19/20)(41/42)(109/110))\simeq0{,}107$.
Jest to dolna granica rezerwy, nie granica samej reszty.
Podział $\R=\sum_{j=1}^K\R_j$, $K\asymp\sqrt X$, daje
$|\R|^2\le K\sum|\R_j|^2$. Wystarczy $\sum|\R_j|^2\ll X^{3/2-\eta}$.
Pomiary $L$ nie są pomiarami $\R_j$.

\section{Dwumian Newtona, zeta i właściwy znak (205--208)}
Rys. 1.3 w notatkach \cite{dupuy} przedstawia
$\binom{20}{10}=2^2\cdot11\cdot13\cdot17\cdot19=184756$ jako ilustrację
formuły iloczynowej (podzielność versus rozmiar); rys. 1.4 przedstawia
$G_\zeta(x)$ do 2000 i znormalizowane fluktuacje do $10^7$. Same wykresy
nie dowodzą RH. Praca Alpögego ogłoszona 20 XI 2026 jest źródłem inspiracji,
nie przesłanką w tym opracowaniu (por. \cite{rub}). Nie ustaliliśmy, czy
wskazany dwumian pełni w formalizacji rolę parametru optymalizacji.

Dla różnic $\Delta^r L_j=\sum_{i=0}^r(-1)^i\binom ri L_{j+i}$ zachodzi
$\sum_i\binom ri^2=\binom{2r}r$. Zatem znormalizowana energia to
\[
 \sigma_r^2=\binom{2r}r^{-1}\operatorname{mean}|\Delta^r L_j|^2.
\]
Wielomian $(1-z)^r$ ma zero rzędu $r$ przy $z=1$ i
$|(1-e^{i\theta})^r|^2=4^r\sin^{2r}(\theta/2)$. Jest to filtr częstotliwości,
nie informacja o zerach dzeta. Test etapu 206 dał m.in. dla $r=3$ brak
osobnego pomiaru, dla $r=4$ wykładniki $0{,}221470/0{,}194560$, dla $r=10$
$0{,}214378/0{,}189515$ (rozłączne/ruchome okna). Mała energia różnic usuwa
również stały błąd systematyczny; nie rozpoznaje znaku.

\subsection{Szereg Dirichleta i wygładzenie}
Dla punktowego $\ell(k)$ z lematu \ref{lem:sign} definiujemy
$F_L(s)=\sum_{k\ge1}\ell(k)k^{-s}$, $\Re s>1$.
Górne sito Selberga w wymiarze 3 zastosowane na $[A,2A)$ przy płytkim
stałopotęgowym progu (mniejszym niż rzeczywisty próg) daje
\begin{equation}\label{eq:upper}
 \sum_{A\le k<2A}|\ell(k)|\ll A/\log^3 A.
\end{equation}
Stąd $F_L(1),F'_L(1)$ są bezwzględnie zbieżne, dla $F''_L(1)$ nie ma
takiej gwarancji. Dla $G_L(x)=\sum_{k\le x}\ell(k)(x/k-1)$ dokładnie
$\int_1^\infty G_L(x)x^{-s-1}dx=F_L(s)/(s(s-1))$ dla $\Re s>1$.
Przeskalowanie $F_L(2s)$ jest dokładne dla $Q=k^2$; dla rzeczywistego
iloczynu $LR=(30k+c)^2-1$ występuje przesunięcie i skala, więc nie jest to
dosłowna tożsamość. Nie znamy przedłużenia meromorficznego $F_L$ ani jego
iloczynu Eulera; osobliwości nie są zerami $\zeta$.

\subsection{Dodatnie jądro Newtona}
\begin{defn}
Dla całkowitych $a\ge1,r\ge3$ i $u>0$ definiujemy
\begin{equation}\label{eq:kernel}
 P_{a,r}(t)=t^a(1-t)^r,\quad
 \B_{a,r}(u)=\sum_{k\ge1}\frac{\ell(k)}{k}P_{a,r}(k^{-u})
 =\sum_{j=0}^r(-1)^j\binom rj F_L(1+(a+j)u).
\end{equation}
\end{defn}
Jądro jest nieujemne na $[0,1]$, ma zero rzędu $a$ przy 0 i $r$ przy 1.
To istotna różnica względem odejmowania kolejnych okien. W szczególności
\begin{tw}[Kryterium znaku]\label{th:kernel}
Dla dowolnych ustalonych $a\ge1,r>2$, jeżeli
\[
 \boxed{\liminf_{u\downarrow0}\frac{\B_{a,r}(u)}{u^2}<0,}
\]
to bliźniaków jest nieskończenie wiele.
\end{tw}
\begin{proof}
Gdyby bliźniaków było skończenie wiele, istniałby $K_0$ taki, że
$\ell(k)\ge0$ dla $k>K_0$. Ponieważ $1-k^{-u}\le u\log k$, mamy
\[
 \B_{a,r}(u)\ge -u^r\sum_{k\le K_0}\frac{|\ell(k)|}{k}(\log k)^r.
\]
Po podzieleniu przez $u^2$ prawa strona dąży do zera. Sprzeczność.
\end{proof}
Jest to elementarna zasada znaku typu Landaua, nie wymaga twierdzenia
Landaua o przedłużeniu szeregu. Ma również wersję: wystarczy, żeby
$\B_{a,r}(u)/u^r$ było nieograniczone od dołu (dla ustalonych parametrów).
Nie wystarczy skończona liczba ujemnych wartości.

\begin{lem}[Bezwarunkowy rząd operatora]
Dla $a\ge1,r>2$ zachodzi $|\B_{a,r}(u)|\ll_{a,r}u^2$ przy $u\downarrow0$.
Wkład zakresu $u\log k\le b<1$ jest $O_{a,r}(u^2b^{r-2}+u^r)$.
\end{lem}
\begin{proof}
Z \eqref{eq:upper}, w bloku $2^j\le k<2^{j+1}$ ważona masa $\sum|\ell|/k$
jest $O(j^{-3})$. Waga jądra jest ograniczona przez
$e^{-a u j\log2}\min(1,(u(j+1)\log2)^r)$. Sumujemy $j\le1/u$ i $j>1/u$.
Pierwsza suma to $O(u^r\sum_{j\le1/u}j^{r-3})=O(u^2)$, druga $O(u^2)$.
Dla odcięcia $b$ zastępujemy $1/u$ przez $b/(u\log2)$ i uwzględniamy
skończony zakres początkowy. Dowód daje ograniczenie rozmiaru, nie znak.
\end{proof}
Przy \emph{hipotetycznej} asymptotyce
$\sum_{k\le A}\ell(k)\sim-\kappa A/\log^3 A$ otrzymujemy przez sumowanie
częściami i podstawienie $v=u\log k$:
\[
 \B_{a,r}(u)\sim-\kappa u^2 I_{a,r},\quad
 I_{a,r}=\int_0^\infty e^{-av}(1-e^{-v})^r\frac{dv}{v^3}>0.
\]
Dla $a=1$ całkowanie przez części dwukrotnie daje
$I_{1,r}=-\frac12\sum_{j=0}^r(-1)^j\binom rj(j+1)^2\log(j+1)$.
\begin{center}
\begin{tabular}{rrrr}
\toprule $r$&$\binom{2r}r$&$I_{1,r}$&mediana $v$ (model)\\\midrule
3&20&0,417972075299&0,294351\\4&70&0,126091300851&0,619178\\
10&184756&0,008867568390&1,693490\\\bottomrule
\end{tabular}
\end{center}
Mniejszy $r$ daje większą całkę, większy $r$ silniej odcina początek.
Wszystkie trzy spełniają kryterium \ref{th:kernel}, jeśli znak zostanie
udowodniony. Wyboru nie można uzasadnić samym podobieństwem do wykresów RH.

\section{Lokalizacja dodatnim wielomianem: audyt etapu 209}\label{sec:cert}
W tym etapie \emph{punktowy} próg $z(k)$ zastosowano dla każdego $k$,
$1\le k\le10^8$. Generator C++ stosuje niezależne maski: pierwszość
przez skreślenia od $p^2$, roughness od pierwszego $k$ z $p^3\le30k+31$.
Dla $p\le5$ wszystkie ramiona są względnie pierwsze z $p$. Tablica 64 na 64
oblicza moment z masek. Ponowne uruchomienie generatora w audycie dało
\textbf{identyczny bajt w bajt} plik \texttt{ell.bin}. SHA-256:
\begin{center}\small\texttt{244fb7963919bb2bb20acf2981d2be9b1a24729192229480135ac6bc91aa624c}\end{center}
Długość: 100000000 bajtów ze znakiem. Suma do $10^7$: $-1555800$;
po $10^7$: $-9242518$; razem $-10798318$. Zakres wartości $[-24,24]$,
wartości parzyste. Współczynniki są dokładne dla podanej definicji, nie
heurystyczne. Niezależna faktoryzacja 1500 cykli (200 początkowych i 1300 wylosowanych) również potwierdziła wszystkie wartości. Właściwa kontrola znaku całej sumy nadal wymaga ogona.

\subsection{Porównanie rzędów (zapisane pomiary)}
Dla $Z_r=\B_{1,r}/(I_{1,r}u^2)$, $u=0{,}1$:
\begin{center}
\begin{tabular}{rrr}
\toprule $r$&$Z_r$ do $10^8$&tylko $k>10^7$\\\midrule
3&$-158{,}636$&$-5{,}947$\\4&$-260{,}051$&$-16{,}171$\\
6&$-329{,}223$&$-40{,}206$\\8&$-317{,}358$&$-59{,}588$\\
10&$-277{,}654$&$-70{,}385$\\12&$-231{,}106$&$-73{,}205$\\
16&$-147{,}473$&$-63{,}696$\\20&$-88{,}378$&$-46{,}871$\\\bottomrule
\end{tabular}
\end{center}
Dla 24 kombinacji $r$ i $u\in\{0{,}05,0{,}1,0{,}2\}$ zapisano ujemność
części skończonej; nie oznacza to znaku pełnej sumy. W celu skupienia masy
na $10^7$--$10^8$ wybrano $u=1/20$, $v_0=u\log 10^{7{,}5}$ i
$a\simeq r/(e^{v_0}-1)-3/v_0$ (zaokrąglenie do dodatniej liczby całkowitej).
Jest to warunek stacjonarności logarytmu modelowego jądra
$e^{-av}(1-e^{-v})^r/v^3$, a nie dopasowywanie znaku.
\begin{center}
\begin{tabular}{rrrrr}
\toprule $r$&$a$&masa modelowa w bloku (\%)&$Z$ skończone&koszt ogona\\\midrule
320&230&74,9843&$-478{,}329$&34390,380\\
640&463&89,6523&$-523{,}934$&12132,056\\
1280&930&97,8894&$-550{,}957$&2240,744\\
2560&1863&99,8872&$-557{,}684$&123,719\\
5120&3730&99,9996&$-558{,}100$&0,515\\\bottomrule
\end{tabular}
\end{center}
Normalizacja $Z=\B_{a,r}/(I_{a,r}u^2)$. Procenty są rozkładem modelowej
masy, nie rozkładem rzeczywistych znaków. W pierwotnym teście kwadratura
oraz część skończona były zmiennoprzecinkowe.

\subsection{Ścisła kontrola całego nieskończonego ogona}
Dla $N\ge1$ niech $v_N=u\log N$, $f(v)=e^{-av}(1-e^{-v})^r$ i
$D=a-r/(e^{v_N}-1)$. Jeśli $D>0$, logarytmiczna pochodna $f$ dla $v\ge v_N$
wynosi $-a+r/(e^v-1)\le-D$. Ponadto $x^{-1}f(u\log x)$ jest malejąca
na $[N,\infty)$. Z kryterium całkowego i $|\ell(k)|\le24$ otrzymujemy
\begin{equation}\label{eq:tail}
 \left|\sum_{k>N}\frac{\ell(k)}k P_{a,r}(k^{-u})\right|
 \le\frac{24}{u}\int_{v_N}^\infty f(v)\,dv
 \le\frac{24 f(v_N)}{uD}.
\end{equation}
To ograniczenie jest bezwarunkowe i obejmuje wszystkie dalsze liczby.
Nie wymaga przewidywania ich pierwszości.

\subsection{Nowy certyfikat z zaokrągleniem skierowanym}
Ustalamy $a=3730$, $r=5120$, $u=1/20$, $N=10^8$. Zamiast numerycznej
całki $I$ stosujemy prosty dodatni czynnik skali $S=10^{2619}$.
Dla $t=k^{-1/20}$ waga jednego wyrazu (bez $\ell(k)$) wynosi
$g(t)=t^{3750}(1-t)^{5120}$. Funkcja ma jeden punkt maksimum
$t_*=3750/8870$. Dla każdego bloku po 100000 cykli odczytujemy dokładne
całkowite sumy dodatnich i ujemnych współczynników $P_b,N_b\ge0$.
Jeżeli $w_b^-\le Sg(t)\le w_b^+$ na tym bloku, wtedy
\[
 S\sum_{k\in b}\frac{\ell(k)}k P_{a,r}(k^{-1/20})
 \le P_bw_b^+-N_bw_b^-.
\]
Przedziały dla $t$ wyznaczamy bez logarytmów: dla $D_0=10^{30}$
znajdujemy całkowite $j$ takie, że $kj^{20}\le D_0^{20}<k(j+1)^{20}$,
stąd $j/D_0\le t\le(j+1)/D_0$. Używamy minimum i maksimum na końcach
bloku, z uwzględnieniem $t_*$, gdy maksimum leży wewnątrz. Potęgi całkowite
liczymy mnożeniami w Decimal z kierunkiem zaokrąglenia ku dołowi lub górze.
Ogon \eqref{eq:tail} ma tu postać $480 f(v_N)/D$; w mianowniku stosujemy
dolny wymierny kres $D=a-r t_N/(1-t_N)$, a w liczniku górny kres.

\begin{center}
\begin{tabular}{lr}
\toprule Składnik (skala $10^{2619}$)&Kres górny\\\midrule
Skończona suma&$-20{,}5147209$\\
Bezwzględny koszt ogona&$0{,}0192070$\\\midrule
Pełna suma&$-20{,}4955139$\\\bottomrule
\end{tabular}
\end{center}
\begin{tw}[Certyfikat dla ustalonych parametrów]
Z danych odtworzonych przez generator i powyższych nierówności:
\[
 \boxed{10^{2619}\B_{3730,5120}(1/20)<-20{,}49.}
\]
\end{tw}
Plik \texttt{certyfikat.py} wykonuje rachunek przedziałowy; wynik zapisano
w \texttt{certyfikat\_wynik.txt}. Jest to certyfikat programistyczny,
nie dowód sprawdzony przez kernel Lean. W szczególności Lean nie przelicza
100 mln współczynników. Ujemność ustalonego $u$ potwierdza tylko istnienie
jakichś ujemnych współczynników (a zatem par). Jeżeli przedłużymy dowolny
skończony zapis przez same zera, wszystkie jego skończone znaki mogą pozostać
ujemne, lecz $\B_{a,r}(u)/u^2\to0$. Żaden wybór rzędu nie usuwa tej przeszkody.

\section{Nowa synteza: B1 wygładzone tym samym jądrem}\label{sec:smooth}
Wcześniejsza praca zawierała dwie poprawne, lecz różne redukcje:
ujemność momentu $L$ i jednostronną resztę Möbiusa B1. Nie ma udowodnionej
tożsamości $L=\R$; to różne wagi arytmetyczne. Można jednak połączyć
\emph{samą redukcję B1} z dodatnim jądrem Newtona, bez takiej tożsamości.
Jest to wspólny cel analityczny całego programu.

Niech $X_j=2^j$, $j\ge1$, oraz
\[
 \tau_j=\T(X_j)/X_j,\qquad r_j=\R(X_j)/X_j,\qquad
 b_j=4C_2(1/2-J_\mu(X_j)/X_j).
\]
Z \eqref{eq:B1exact}, \eqref{eq:reserve} mamy
$\tau_j=b_j+r_j+\epsilon_j$, $\epsilon_j\to0$, $b_j\ge\delta_0+o(1)$.
Sito górne dla bliźniaków zapewnia $\tau_j=O(1)$, a więc także $r_j=O(1)$.
Wybieramy te same dodatnie jądra co wyżej, lecz na osi indeksu dyadycznego:
\begin{equation}\label{eq:Wj}
 W_j(u)=j^{-3}e^{-au j}(1-e^{-u j})^r,\qquad
 W(u)=\sum_{j\ge1}W_j(u),\quad a\ge1,\ r>2.
\end{equation}
\begin{lem}
$W(u)\sim I_{a,r}u^2$ i dla każdego ustalonego $J$,
$\sum_{j\le J}W_j(u)=O_{J,r}(u^r)=o(W(u))$.
Jeżeli $e_j\to0$, to $\sum_j e_jW_j(u)=o(W(u))$.
\end{lem}
\begin{proof}
Podstawienie $v=uj$ daje $W(u)=u^2[u\sum_{j\ge1}f(uj)/(uj)^3]$;
wyrażenie w nawiasie jest sumą Riemanna zbieżną do $I_{a,r}$.
Zbieżność przy zerze wynika z $r>2$, przy nieskończoności z $a>0$.
Dla ustalonego $j$ mamy $1-e^{-uj}\le uj$. Ostatni wniosek wynika przez
podział na skończony początek i ogon, gdzie $|e_j|$ jest dowolnie małe.
\end{proof}
\begin{tw}[Główny lemat redukcyjny: B1--Newton]\label{th:B1N}
Wystarczy udowodnić dla jednego ustalonego $a\ge1,r>2$:
\begin{equation}\label{eq:B1N}
 \boxed{\limsup_{u\downarrow0}
 \frac{\sum_{j\ge1}W_j(u)\R(2^j)/2^j}{\sum_{j\ge1}W_j(u)}>-\delta_0,
 \qquad \delta_0=C_2/10.}
\end{equation}
Wtedy istnieje nieskończenie wiele bliźniaków.
\end{tw}
\begin{proof}
Pomnożenie $\tau_j\ge\delta_0+r_j+o(1)$ przez nieujemne $W_j$ i sumowanie
po $j$ daje
\[
 \frac{\sum_j W_j\tau_j}{W(u)}\ge\delta_0+
 \frac{\sum_j W_j r_j}{W(u)}+o(1).
\]
Jeżeli zachodzi \eqref{eq:B1N}, prawa strona jest dodatnia i oddzielona
od zera wzdłuż ciągu $u\downarrow0$. Przy skończenie wielu bliźniakach
$\tau_j=0$ poza skończonym zbiorem, a ich średnia ważona dąży do zera.
Sprzeczność.
\end{proof}
Ta wersja nie wymaga małości każdej lokalnej rozbieżności, tylko jednostronnej
średniej podpisanej. Jest kompatybilna z filtrem Newtona, z CRT v10
($\Delta_e$), z naturalnym dzieleniem bloków na okna Czybyszewa i
z pierwotnym B1. Można zwiększyć $\delta_0$ do $\delta_*$ po uwzględnieniu
kolejnych kwadratów. Nie twierdzimy, że \eqref{eq:B1N} wynika z pomiaru
$\B_{3730,5120}(1/20)$: ten ostatni dotyczy $\ell(k)$, nie $r_j$.
Nie twierdzimy też, że punktowe B1 na dowolnym rzadkim ciągu implikuje
\eqref{eq:B1N}; oba są osobnymi warunkami wystarczającymi. B1 wymusza ilościową dolną granicę rzędu $X/\log^2 X$, podczas gdy kryterium znaku momentu $L$ może dać nieskończoność przy słabszej masie; nie są równoważne.

\section{Plan dalszych badań: wyłącznie właściwa bariera}
Przedmiotem dalszej analizy powinien być wyraz $\R(X)$ z \eqref{eq:R}
(z wagami, ruchomą granicą i centrowaniem), a nie ponowne liczenie zwykłych
bliźniaków. Docelowy warunek to B1 albo jego wersja \eqref{eq:B1N}.
Wspólne dzielniki $e\mid n+2$ oznaczają dokładnie prawe klasy CRT. Z wcześniejszych
etapów należy zachować metody Kloostermana, biliniowe przełączanie
czynników, dyspersję oraz signed energy jako potencjalne narzędzia; nie
uznano żadnej z nich za pełne oszacowanie dla wszystkich zakresów.
Również duże czynniki środka $C_m$ i równania norm są narzędziami
pomocniczymi, nie koniecznymi elementami redukcji.

Trzy sprawdzalne zadania: (1) rozkład $\R$ na bloki czynników z dokładnym
zachowaniem znaków i granic; (2) oszacowanie jednostronnej sumy bloków,
których nie pokrywa BV ani klasyczna dyspersja; (3) zbadanie dodatniego
jądra \eqref{eq:Wj} bez przenoszenia błędu z $L$ na $\R$ bez dowodu.
Jeżeli połączony lemat zostanie udowodniony, twierdzenie \ref{th:B1N}
kończy dowód. Bez niego pozostaje bariera parzystości. Oś $1/4$, zmiana
parametrów i dalsze testy nie stanowią samodzielnego przełamania bariery.

\appendix
\section{Pakiet reprodukcji i Lean}
W archiwum: niniejsze źródło LaTeX i PDF, \texttt{AudytRK.lean},
\texttt{certyfikat.py}, \texttt{audyt\_danych.py}, raporty, dane i kody z etapów 193,
197, 199--203 i 209 oraz wcześniejszy tekst 140--144 (źródło odniesienia).
Zasadniczy generator i certyfikat uruchamia się następująco (Linux/WSL):
\begin{verbatim}
g++ -O3 -fopenmp dane/etap209/test.cpp -o rk209
./rk209 100000000 ell.bin 2> bloki.csv
python certyfikat.py ell.bin
python audyt_danych.py dane
sha256sum ell.bin
\end{verbatim}
Wymagany Python z NumPy; poprzednie analizy używają także SciPy i pandas.
W projekcie Lean z Mathlib: \texttt{lake env lean AudytRK.lean}.
Plik Lean nie był uruchomiony w środowisku przygotowania (brak Lean/Mathlib).
Nie dodano deklaracji nowych aksjomatów ani świadomych luk \texttt{sorry}.
Przesłanki analityczne są argumentami twierdzeń, a nie udowodnionymi lematami.
Do sprawdzenia: składnia, użyte nazwy lematów bibliotecznych i komunikaty
\texttt{\#print axioms}. Rdzeń Lean obejmuje algebraiczny fundament,
momenty, rzeczywistą implikację $\ell(k)<0\Rightarrow$ pełny pierwszy tor,
obliczenia wymierne, teleskop, dodatnie jądra i logiczną implikację
z odpowiednio silnego znaku. Formalizacja pełnej analizy BV, szeregu
Möbiusa, asymptotycznej sumy Riemanna i obliczeń zewnętrznych nie została
wykonana. Plik nie jest formalnym dowodem hipotezy bliźniaczej.

\begin{thebibliography}{9}
\bibitem{mv} M. R. Murty, A. Vatwani, \emph{Twin primes and the parity problem},
Journal of Number Theory 180 (2017), 643--659, zwł. lematy 3.4, 4.4.
\url{https://mrammurty-ai.github.io/TwinPrimes-Parity.pdf}.
\bibitem{taosieve} T. Tao, \emph{254A, Notes 4: Some sieve theory}, 2015.
\url{https://terrytao.wordpress.com/2015/01/21/254a-notes-4-some-sieve-theory/}.
\bibitem{tt} T. Tao, J. Teräväinen, \emph{The structure of correlations of
multiplicative functions at almost all scales, with applications to the
Chowla and Elliott conjectures}, arXiv:1809.02518.
\url{https://arxiv.org/abs/1809.02518}.
\bibitem{dupuy} Notatki o ogłoszonej formalizacji Alpögego--Claude,
wersja dostępna u T. Dupuy, 21 XI 2026, rozdz. 1, rys. 1.3--1.4.
\url{https://tdupu.github.io/riemann-public/riemann-book.pdf}.
\bibitem{rub} M. Rubinstein, \emph{The Alpöge--Claude proof of the Riemann
hypothesis explained for the Riemann zeta function}, working draft v2,
22 XI 2026. \url{https://math.uwaterloo.ca/~mrubinst/alpoge_claude_grh/zeta-paper_v2.pdf}.
\end{thebibliography}
\end{document}
